\documentclass[12pt]{article}
\usepackage[a4paper,margin=1in]{geometry}
\usepackage[utf8]{inputenc}
\usepackage[T1]{fontenc}
\usepackage{amsmath,amssymb}
\usepackage{textcomp}
\usepackage{parskip}
\begin{document}
\begin{center}
{\Large \textbf{2004 STEP I}}
\end{center}

\subsection*{Page 1}
\begin{flushleft}
STEP I, 2004 2\par
\bigskip\textbf{Section A: Pure Mathematics}\par
\[1 (i) Express 3 + 2\sqrt\{5\} in the form a + b\sqrt\{5\} where a and b are integers.\]
\textbf{(ii)} Find the positive integers c and d such that 3 99 70\sqrt2 = c d\sqrt2 .\par
p - -\par
\textbf{(iii)} Find the two real solutions of x6 198x3 + 1 = 0 .\par
-\par
\bigskip\textbf{Question 2.} The square bracket notation [x] means the greatest integer less than or equal to x . For\par
\[example, [\pi] = 3 , [\sqrt\{24\} ] = 4 and [5] = 5 .\]
\textbf{(i)} Sketch the graph of y = [x] and show that\par
p\par
a a 1\par
-\par
\[[x] dx = \sqrt\{r\}\]
\[\int\]
\[0 p Xr=0\]
when a is a positive integer.\par
a\par
[[[x]]] a\par
\textbf{(ii)} Show that 2 dx = 2 1 when a is a positive integer.\par
Z -\par
a\par
[[[x]]]\par
\textbf{(iii)} Determine an expression for 2 dx when a is positive but not an integer.\par
\[\int\]
\bigskip\textbf{Question 3.} (i) Show that x 3 is a factor of\par
-\par
3 2 2 2 2\par
x 5x + 2x y + xy 8xy 3y + 6x + 6y . ( )\par
- - - ∗\par
Express ( ) in the form (x 3)(x + ay + b)(x + cy + d) where a, b, c and d are integers\par
∗ -\par
to be determined.\par
\textbf{(ii)} Factorise 6y3 y2 21y + 2x2 + 12x 4xy + x2y 5xy2 + 10 into three linear factors.\par
- - - -\par
\end{flushleft}
\subsection*{Page 2}
\begin{flushleft}
STEP I, 2004 3\par
4 Differentiate sec t with respect to t.\par
\[2 1 \sqrt\{3\} 2 \pi\]
\textbf{(i)} Use the substitution x = sec t to show that dx = - + .\par
\[Z \sqrt\{2\} x^3\sqrt\{x^2\} 1 8 24\]
-\par
\textbf{(ii)} Determine dx .\par
Z (x + 2) (x + 1)(x + 3)\par
p\par
\textbf{(iii)} Determine dx .\par
\[Z (x + 2)\sqrt\{x^2\} + 4x 5\]
-\par
\bigskip\textbf{Question 5.} The positive integers can be split into five distinct arithmetic progressions, as shown:\par
A : 1, 6, 11, 16, ...\par
B : 2, 7, 12, 17, ...\par
C : 3, 8, 13, 18, ...\par
D : 4, 9, 14, 19, ...\par
E : 5, 10, 15, 20, ...\par
Write down an expression for the value of the general term in each of the five progressions.\par
Hence prove that the sum of any term in B and any term in C is a term in E.\par
Prove also that the square of every term in B is a term in D. State and prove a similar claim\par
about the square of every term in C.\par
\textbf{(i)} Prove that there are no positive integers x and y such that\par
\[x + 5y = 243 723 .\]
\textbf{(ii)} Prove also that there are no positive integers x and y such that\par
4 4\par
\[x + 2y = 26 081 974 .\]
\end{flushleft}
\subsection*{Page 3}
\begin{flushleft}
STEP I, 2004 4\par
\bigskip\textbf{Question 6.} The three points A, B and C have coordinates (p 1 , q 1), (p 2 , q 2) and (p 3 , q 3) , respectively.\par
Find the point of intersection of the line joining A to the midpoint of BC, and the line\par
joining B to the midpoint of AC. Verify that this point lies on the line joining C to the\par
midpoint of AB.\par
The point H has coordinates (p 1 + p 2 + p 3 , q 1 + q 2 + q 3) . Show that if the line AH intersects\par
the line BC at right angles, then p2 + q2 = p2 + q2 , and write down a similar result if the line\par
2 2 3 3\par
BH intersects the line AC at right angles.\par
Deduce that if AH is perpendicular to BC and also BH is perpendicular to AC, then CH is\par
perpendicular to AB.\par
\bigskip\textbf{Question 7.} (i) The function f (x) is defined for x < 1 by\par
| | 5\par
\infty n\par
\[f (x) = anx ,\]
\[nX=0\]
\[where a 0 = 2, a 1 = 7 and an = 7an 1 10an 2 for n > 2 .\]
- - -\par
1 1\par
\[Simplify f (x) 7xf (x) + 10x2f (x) , and hence show that f (x) = + .\]
- 1 2x 1 5x\par
n n - -\par
\[Hence show that an = 2 + 5 .\]
\textbf{(ii)} The function g (x) is defined for x < 1 by\par
| | 3\par
\infty n\par
\[g (x) = bnx ,\]
\[nX=0\]
\[where b 0 = 5 , b 1 = 10 , b 2 = 40 , b 3 = 100 and bn = pbn 1 + qbn 2 for n > 2 . Obtain\]
- -\par
an expression for g (x) as the sum of two algebraic fractions and determine bn in terms\par
of n .\par
\bigskip\textbf{Question 8.} A sequence t 0, t 1, t 2, ... is said to be strictly increasing if tn+1 > tn for all n > 0 .\par
\textbf{(i)} The terms of the sequence x 0 , x 1 , x 2 , . . . satisfy\par
\[x^2 n + 6\]
\[xn+1 =\]
\[for n > 0 . Prove that if x 0 > 3 then the sequence is strictly increasing.\]
\textbf{(ii)} The terms of the sequence y 0 , y 1 , y 2 , . . . satisfy\par
\[yn+1 = 5\]
- yn\par
\[for n > 0 . Prove that if 2 < y 0 < 3 then the sequence is strictly increasing but that\]
\[yn < 3 for all n .\]
\end{flushleft}
\subsection*{Page 4}
\begin{flushleft}
STEP I, 2004 5\par
\bigskip\textbf{Section B: Mechanics}\par
\bigskip\textbf{Question 9.} A particle is projected over level ground with a speed u at an angle θ above the horizontal.\par
Derive an expression for the greatest height of the particle in terms of u, θ and g.\par
A particle is projected from the floor of a horizontal tunnel of height 9 d. Point P is 1 d\par
10 2\par
metres vertically and d metres horizontally along the tunnel from the point of projection. The\par
particle passes through point P and lands inside the tunnel without hitting the roof. Show\par
that\par
\[arctan 3 < θ < arctan 3 .\]
\bigskip\textbf{Question 10.} A particle is travelling in a straight line. It accelerates from its initial velocity u to velocity v,\par
where v > u > 0 , travelling a distance d 1 with uniform acceleration of magnitude 3a . It then\par
| |\par
comes to rest after travelling a further distance d 2 with uniform deceleration of magnitude a .\par
Show that\par
\textbf{(i)} if u > 0 then 3d 1 < d 2 ;\par
\textbf{(ii)} if u < 0 then d 2 < 3d 1 < 2d 2 .\par
Show also that the average speed of the particle (that is, the total distance travelled divided\par
\[by the total time) is greater in the case u > 0 than in the case u < 0 .\]
Note: In this question d 1 and d 2 are distances travelled by the particle which are not the\par
same, in the second case, as displacements from the starting point.\par
\bigskip\textbf{Question 11.} Two uniform ladders AB and BC of equal length are hinged smoothly at B. The weight of\par
AB is W and the weight of BC is 4W . The ladders stand on rough horizontal ground with\par
\\par
\[ABC = 60 . The coefficient of friction between each ladder and the ground is µ.\]
\textasciicircum{}\circ\par
A decorator of weight 7W begins to climb the ladder AB slowly. When she has climbed up 1\par
of the ladder, one of the ladders slips. Which ladder slips, and what is the value of µ?\par
\end{flushleft}
\subsection*{Page 5}
\begin{flushleft}
STEP I, 2004 6\par
\bigskip\textbf{Section C: Probability and Statistics}\par
\bigskip\textbf{Question 12.} In a certain factory, microchips are made by two machines. Machine A makes a proportion λ\par
\[of the chips, where 0 < λ < 1, and machine B makes the rest. A proportion p of the chips\]
made by machine A are perfect, and a proportion q of those made by machine B are perfect,\par
where 0 < p < 1 and 0 < q < 1. The chips are sorted into two groups: group 1 contains those\par
that are perfect and group 2 contains those that are imperfect.\par
In a large random sample taken from group 1, it is found that were made by machine A.\par
Show that λ can estimated as\par
2q\par
.\par
3p + 2q\par
Subsequently, it is discovered that the sorting process is faulty: there is a probability of that\par
a perfect chip is assigned to group 2 and a probability of that an imperfect chip is assigned\par
to group 1. Taking into account this additional information, obtain a new estimate of λ .\par
13 (i) Three real numbers are drawn independently from the continuous rectangular distribu-\par
tion on [0, 1] . The random variable X is the maximum of the three numbers. Show that\par
the probability that X 6 0.8 is 0.512 , and calculate the expectation of X.\par
\textbf{(ii)} N real numbers are drawn independently from a continuous rectangular distribution\par
on [0, a] . The random variable X is the maximum of the N numbers. A hypothesis\par
test with a significance level of 5\% is carried out using the value, x, of X. The null\par
\[hypothesis is that a = 1 and the alternative hypothesis is that a < 1 . The form of the\]
\[test is such that H 0 is rejected if x < c , for some chosen number c .\]
Using the approximation 210 103 , determine the smallest integer value of N such that\par
\approx\par
if x 6 0.8 the null hypothesis will be rejected.\par
With this value of N , write down the probability that the null hypothesis is rejected if\par
\[a = 0.8 , and find the probability that the null hypothesis is rejected if a = 0.9 .\]
\end{flushleft}
\subsection*{Page 6}
\begin{flushleft}
STEP I, 2004 7\par
\bigskip\textbf{Question 14.} Three pirates are sharing out the contents of a treasure chest containing n gold coins and 2\par
lead coins. The first pirate takes out coins one at a time until he takes out one of the lead\par
coins. The second pirate then takes out coins one at a time until she draws the second lead\par
coin. The third pirate takes out all the gold coins remaining in the chest.\par
Find:\par
\textbf{(i)} the probability that the first pirate will have some gold coins;\par
\textbf{(ii)} the probability that the second pirate will have some gold coins;\par
\textbf{(iii)} the probability that all three pirates will have some gold coins.\par
\end{flushleft}
\end{document}
